\clearpage\selectlanguage{ngerman}\fancyhead[L]{\small ALLGEMEINE UNTERSUCHUNGEN}
\SourceLeaf{134}
\begin{center}
ALLGEMEINE UNTERSUCHUNGEN

\medskip
ÜBER DIE UNENDLICHE REIHE
\[
\begin{aligned}&1+\frac{\alpha\beta}{1\mathbin{.}\gamma}x\\
&{}+\frac{\alpha(\alpha+1)\beta(\beta+1)}{1\mathbin{.}2\mathbin{.}\gamma(\gamma+1)}xx\\
&{}+\frac{\alpha(\alpha+1)(\alpha+2)\beta(\beta+1)(\beta+2)}{1\mathbin{.}2\mathbin{.}3\mathbin{.}\gamma(\gamma+1)(\gamma+2)}x^3+\mathrm{usw.}
\end{aligned}
\]
ERSTER TEIL.

\medskip
\rule{0.13\linewidth}{0.4pt}

\medskip
\textit{EINLEITUNG.}

\end{center}
\SixthArticle{1}
Die Reihe, deren Untersuchung wir in dieser Abhandlung unternehmen, kann als Funktion von vier Größen $\alpha,\beta,\gamma,x$ betrachtet werden, die wir ihre \textit{Elemente} nennen werden. Wir unterscheiden sie der Reihe nach als erstes Element $\alpha$, zweites $\beta$, drittes $\gamma$ und viertes $x$. Offenbar lässt sich das erste Element mit dem zweiten vertauschen. Bezeichnen wir also der Kürze halber unsere Reihe mit $F(\alpha,\beta,\gamma,x)$, so gilt $F(\beta,\alpha,\gamma,x)=F(\alpha,\beta,\gamma,x)$.

\SixthArticle{2}
Weist man den Elementen $\alpha,\beta,\gamma$ bestimmte Werte zu, so wird unsere Reihe zu einer Funktion der einzigen Variablen $x$. Diese bricht offenbar nach dem $1-\alpha^{\mathrm{ten}}$ oder $1-\beta^{\mathrm{ten}}$ Glied ab, wenn $\alpha-1$ oder $\beta-1$ eine negative ganze Zahl ist; in den übrigen Fällen setzt sie sich dagegen ins Unendliche fort. Im ersten Fall stellt die Reihe eine rationale algebraische Funktion dar, im zweiten dagegen meistens eine transzendente Funktion. Das dritte Element $\gamma$ darf weder eine negative ganze Zahl noch $=0$ sein, damit wir nicht zu unendlich großen Gliedern gelangen.
\SourceLeaf{135}
\SixthArticle{3}
Die Koeffizienten der Potenzen $x^m$, $x^{m+1}$ in unserer Reihe verhalten sich wie
\[
1+\frac{\gamma+1}{m}+\frac{\gamma}{mm}:1+\frac{\alpha+\beta}{m}+\frac{\alpha\beta}{mm}
\]
und ihr Verhältnis nähert sich daher umso mehr der Einheit, je größer $m$ angenommen wird. Wird also auch dem vierten Element $x$ ein bestimmter Wert zugewiesen, so hängt die Konvergenz oder Divergenz von dessen Beschaffenheit ab. Wird nämlich $x$ ein reeller Wert, positiv oder negativ, kleiner als die Einheit zugewiesen, so wird die Reihe sicher, wenn nicht gleich von Anfang an, so doch nach einem gewissen Abschnitt konvergieren und zu einer endlichen, vollständig bestimmten Summe führen. Dasselbe tritt bei einem imaginären Wert von $x$ der Form $a+b\BareRoot{-1}$ ein, wenn $aa+bb<1$ ist. Dagegen wird die Reihe für einen reellen Wert von $x$, der größer als die Einheit ist, oder für einen imaginären Wert der Form $a+b\BareRoot{-1}$, wenn $aa+bb>1$ ist, notwendig divergieren, wenn nicht sofort, so doch nach einem gewissen Abschnitt, sodass von ihrer \textit{Summe} nicht die Rede sein kann. Schließlich hängt für den Wert $x=1$ (oder allgemeiner für einen Wert der Form $a+b\BareRoot{-1}$, wenn $aa+bb=1$ ist) die Konvergenz oder Divergenz der Reihe von der Beschaffenheit von $\alpha,\beta,\gamma$ selbst ab. Darüber und insbesondere über die Summe der Reihe für $x=1$ werden wir im dritten Abschnitt sprechen.

Es ist daher klar, dass die Untersuchung, soweit unsere Funktion als Summe einer Reihe definiert ist, ihrer Natur nach auf diejenigen Fälle beschränkt ist, in denen die Reihe tatsächlich konvergiert. Es ist also eine unsinnige Frage, welchen Wert die Reihe für einen Wert von $x$ hat, der größer als die Einheit ist. Weiter unten, vom vierten Abschnitt an, werden wir unsere Funktion jedoch auf ein höheres Prinzip gründen, das die allgemeinste Anwendung zulässt.

\SixthArticle{4}
Die Differentiation unserer Reihe, bei der allein das vierte Element $x$ als variabel betrachtet wird, führt auf eine ähnliche Funktion, da offenbar gilt
\[
\frac{\mathrm{d}F(\alpha,\beta,\gamma,x)}{\mathrm{d}x}
=\frac{\alpha\beta}{\gamma}F(\alpha+1,\beta+1,\gamma+1,x)
\]
Dasselbe gilt für wiederholte Differentiationen.

\SixthArticle{5}
Es wird sich lohnen, hier einige Funktionen anzuführen, die sich auf unsere Reihe zurückführen lassen und in der gesamten Analysis sehr häufig gebraucht werden.
\SourceLeaf{136}
\[
\begin{array}{rl}
\mathrm{I.}&(t+u)^n=t^nF(-n,\beta,\beta,-\frac{u}{t})
\end{array}
\]
wobei das Element $\beta$ beliebig ist.
\[
\begin{array}{rl}
\mathrm{II.}&(t+u)^n+(t-u)^n=2t^nF(-\frac{1}{2}n,-\frac{1}{2}n+\frac{1}{2},\frac{1}{2},\frac{uu}{tt})
\end{array}
\]
\[
\begin{array}{rl}
\mathrm{III.}&(t+u)^n+t^n=2t^nF(-n,\omega,2\omega,-\frac{u}{t})
\end{array}
\]
wobei $\omega$ eine unendlich kleine Größe bezeichnet.
\[
\begin{array}{rl}
\mathrm{IV.}&(t+u)^n-(t-u)^n=2nt^{n-1}uF(-\frac{1}{2}n+\frac{1}{2},-\frac{1}{2}n+1,\frac{3}{2},\frac{uu}{tt})
\end{array}
\]
\[
\begin{array}{rl}
\mathrm{V.}&(t+u)^n-t^n=nt^{n-1}uF(1-n,1.\ 2,-\frac{u}{t})
\end{array}
\]
\[
\begin{array}{rl}
\mathrm{VI.}&\log(1+t)=tF(1,1,2,-t)
\end{array}
\]
\[
\begin{array}{rl}
\mathrm{VII.}&\log\frac{1+t}{1-t}=2tF(\frac{1}{2},1,\frac{3}{2},tt)
\end{array}
\]
\[
\begin{array}{rl}
\mathrm{VIII.}&e^t=F(1,k,1,\frac{t}{k})=1+tF(1,k,2,\frac{t}{k})=1+t+\frac{1}{2}ttF(1,k,3,\frac{t}{k})\ \mbox{usw.}
\end{array}
\]
wobei $e$ die Basis der natürlichen Logarithmen und $k$ eine unendlich große Zahl bezeichnet.
\[
\begin{array}{rl}
\mathrm{IX.}&e^t+e^{-t}=2F(k,k',\frac{1}{2},\frac{tt}{4kk'})
\end{array}
\]
wobei $k$, $k'$ unendlich große Zahlen bezeichnen.
\[
\begin{array}{rl}
\mathrm{X.}&e^t-e^{-t}=2tF(k,k',\frac{3}{2},\frac{tt}{4kk'})
\end{array}
\]
\[
\begin{array}{rl}
\mathrm{XI.}&\sin t=tF(k,k',\frac{3}{2},-\frac{tt}{4kk'})
\end{array}
\]
\[
\begin{array}{rl}
\mathrm{XII.}&\cos t=F(k,k',\frac{1}{2},-\frac{tt}{4kk'})
\end{array}
\]
\[
\begin{array}{rl}
\mathrm{XIII.}&t=\sin t.\ F(\frac{1}{2},\frac{1}{2},\frac{3}{2},\sin t^2)
\end{array}
\]
\[
\begin{array}{rl}
\mathrm{XIV.}&t=\sin t.\ \cos t.\ F(1,1,\frac{3}{2},\sin t^2)
\end{array}
\]
\[
\begin{array}{rl}
\mathrm{XV.}&t=\mathrm{tang}\,t.\ F(\frac{1}{2},1,\frac{3}{2},-\mathrm{tang}\,t^2)
\end{array}
\]
\[
\begin{array}{rl}
\mathrm{XVI.}&\sin nt=n\sin t.\ F(\frac{1}{2}n+\frac{1}{2},-\frac{1}{2}n+\frac{1}{2},\frac{3}{2},\sin t^2)
\end{array}
\]
\[
\begin{array}{rl}
\mathrm{XVII.}&\sin nt=n\sin t.\ \cos t.\ F(\frac{1}{2}n+1,-\frac{1}{2}n+1,\frac{3}{2},\sin t^2)
\end{array}
\]
\[
\begin{array}{rl}
\mathrm{XVIII.}&\sin nt=n\sin t.\ \cos t^{n-1}F(-\frac{1}{2}n+1,-\frac{1}{2}n+\frac{1}{2},\frac{3}{2},-\mathrm{tang}\,t^2)
\end{array}
\]
\[
\begin{array}{rl}
\mathrm{XIX.}&\sin nt=n\sin t.\ \cos t^{-n-1}F(\frac{1}{2}n+1,\frac{1}{2}n+\frac{1}{2},\frac{3}{2},-\mathrm{tang}\,t^2)
\end{array}
\]
\[
\begin{array}{rl}
\mathrm{XX.}&\cos nt=F(\frac{1}{2}n,-\frac{1}{2}n,\frac{1}{2},\sin t^2)
\end{array}
\]
\[
\begin{array}{rl}
\mathrm{XXI.}&\cos nt=\cos t.\ F(\frac{1}{2}n+\frac{1}{2},-\frac{1}{2}n+\frac{1}{2},\frac{1}{2},\sin t^2)
\end{array}
\]
\[
\begin{array}{rl}
\mathrm{XXII.}&\cos nt=\cos t^nF(-\frac{1}{2}n,-\frac{1}{2}n+\frac{1}{2},\frac{1}{2},-\mathrm{tang}\,t^2)
\end{array}
\]
\[
\begin{array}{rl}
\mathrm{XXIII.}&\cos nt=\cos t^{-n}F(\frac{1}{2}n+\frac{1}{2},\frac{1}{2}n,\frac{1}{2},-\mathrm{tang}\,t^2)
\end{array}
\]
